% NOETHER PAPER 2 — KOREAN TRANSLATION-PRODUCER DRAFT — UNCHECKED
% Unit: T01-U02; current-authority whole lines 483--496
% Source slice: 2,655 LF UTF-8 bytes; SHA-256 D9660703F6FBF0A3106AEF04FAB4777FCDD07121D402CE0C728E94CC713D7012
% Pointer: NOETH-DE-AUTH-v004-20260804; SHA-256 A1C62FDACAA34DFC1B806DC18258F2E732539F7AE9D85AA4BD9E1067B8749D9F
% This is producer translation only: no source/Korean/formula review, compilation, or rendering.

\begin{center}
\textbf{서론.}
\end{center}

삼원 4차 형식의 형식계는 \emph{Gordan}, \emph{Maisano}, \emph{Pascal}의 연구에서 다루어졌다\srcfn{**)}{\emph{P. Gordan}, Über das volle Formensystem der ternären biquadratischen Form
\[
f=x_1^3x_2+x_2^3x_3+x_3^3x_1.
\]
(Math. Annalen, Bd. XVII (1880), S. 217--233.)
\emph{G. Maisano}, 1) \emph{Sistemi completi dei primi cinque gradi della forma ternaria biquadratica e degl' invarianti, covarianti e contravarianti di sesto grado}. (Giorn. di Battaglini XIX.)
2) \emph{Sui covarianti indipendenti di $6^\circ$ grado nei coefficienti della forma biquadratica ternaria}. (Rend. Circ. Mat. di Palermo I, 1887.)
\emph{E. Pascal}, \emph{Contributo alla teoria della forma ternaria biquadratica e delle sue varie decomposizioni in fattori}. (Memoria premiata dalla R. Accademia delle scienze fisiche e matematiche di Napoli. 1905.)}. \emph{Gordan}은 이항 영역의 형식계에 사용했던 것과 비슷한 원리를 바탕으로, 특수한 자기동형 형식 (f=x_1^3x_2+x_2^3x_3+x_3^3x_1)의 54개 구성식으로 이루어진 완전한 형식계를 세웠다.

\emph{Maisano}는 일반 4차 형식에 대해 계수차수 5까지의 형식들\srcfn{***)}{여기서 \glqq{}Ordnung\grqq{}은 계수에 관한 차수를, \glqq{}Grad\grqq{}는 변수에 관한 차수를 뜻한다.}과 더 높은 계수차수의 몇몇 불변식, 공변식, 반공변식을 제시하였다. 그 방법은 \emph{Gordan}이 Math. Annalen 제1권에서 삼원 3차 형식에 적용한 것이다. \emph{Pascal}은 \emph{Maisano}의 결과를 이용하여 주로 4차 형식의 인수분해 문제를 다루었다\srcfn{*)}{두 번째 짧은 논문에서 \emph{Maisano}는 계수차수 6, 변수차수 6인 세 공변식의 선형 종속성을 보이려 하였다. 완전하지 못한 이 증명과 달리 \emph{Pascal}은 계수차수 6인 세 공변식의 선형 독립성과 계수차수 9인 세 불변식의 선형 독립성을 증명했다고 여겼다. 그러나 실제로 한편의 세 공변식 사이에, 다른 한편의 세 불변식 사이에 각각 선형 관계가 존재한다는 사실은 이 논문의 공식 (13), § 11과 공식 (22), § 17 주석에서 확인되며, 그곳에 이 관계들이 명시적으로 주어져 있다. \emph{Pascal}의 통지에 따르면 이 모순은 그의 논문 46쪽에 있는 반공변식 $p$(여기서는 $\varrho$라 부른다)의 식에 생긴 수치 오류로 설명된다.}.

